\emph{Conventions.} The colour factor is the $\mathfrak{su}(3)$ summand of Theorem \claimref{standard-model.gauge-algebra-minimal-massive-composite}, realised in the colour block as the bivectors of $V$ commuting with $J$ and $\mathrm{Re}\,\Omega$ (the check solves for them: eight). Lattice weights and Haar measure are those of Fact \claimref{synthesis.imported-wilson-lattice-gauge-theory-and-strong-coupling-confinement}(i), at $\beta = 0$. The frozen-matter argument follows Step 4 of the proof of \claimref{synthesis.non-abelian-gauge-fixing-has-gribov-copies-and-the-weak-rotor-lattice-with-doublets-has-a-perimeter-law}, with the doublet replaced by a triplet.

\emph{Step 1 (moments; computed).} For a compact group the average of $U^{\otimes k}$ is the orthogonal projector onto the invariant tensors of $V^{\otimes k}$ (invariance of the measure gives $P^{2} = P$ with image the invariants, inversion invariance gives $P^{\mathsf T} = P$). The finite group $G$ of $J$-commuting signed permutations of $e_{3}, \dots, e_{8}$ that fix $\mathrm{Re}\,\Omega$ (order $96$) lies in $SU(3)$, so the $SU(3)$-invariants are the $G$-invariants annihilated by the eight $\mathfrak{su}(3)$ derivations ($SU(3)$ is connected). The check computes them in degree two (two: the metric and the form of $J$) and degree four (twelve, inside the twenty $G$-invariants), and reads off $\int x^{2} = \tfrac16$, $\int x^{4} = \tfrac1{16}$, $\int U_{ab}x = \tfrac16\Pi_{ab}$, $\int U_{ab}x^{3} = \tfrac1{16}\Pi_{ab}$ as entries of the projectors $B(B^{\mathsf T}B)^{-1}B^{\mathsf T}$. The average of $x^{4}$ over $G$ is $\tfrac16$ (sixteen of the $96$ elements send $e_{3}$ to $\pm e_{3}$), so degree four is not shared. (As a cross-check by another route: $Ue_{3}$ is uniform on the unit sphere of $V$, on which the fourth moment of a coordinate is $3/(6 \cdot 8) = \tfrac1{16}$.) This is (i).

\emph{Step 2 (unitary gauge; proved).} $SU(3)$ acts transitively on the unit sphere of $V$, so for each site choose $g_{x} \in SU(3)$ with $g_{x}e_{3} = \varphi_{x}$; then $\varphi_{x}\cdot U(\varphi_{y}) = e_{3}\cdot(g_{x}^{-1}Ug_{y})(e_{3})$ (the check confirms this on a sample of elements of $G$). The change of variables $U_{\ell} \mapsto g_{x}^{-1}U_{\ell}g_{y}$ preserves the Haar measure and $\chi(U_{C})$ for closed $C$, so the $\varphi$ integrals factor out and each link carries the independent weight $e^{\kappa x(U)}$.

\emph{Step 3 (the mean link; proved).} Let $M = \int Ue^{\kappa x}dU/\int e^{\kappa x}dU$. The stabiliser $H$ of $e_{3}$ in $SU(3)$ fixes $e_{4} = j(e_{3})$ and acts on $\mathrm{span}(e_{5}, \dots, e_{8})$ as $SU(2)$, with no fixed vector there. Since $x(hUh') = x(U)$ for $h, h' \in H$, $M = hMh'$: with $h' = 1$ the image of $M$ lies in the $H$-fixed vectors $\mathrm{span}(e_{3}, e_{4})$; with $h = 1$, $M(h'v - v) = 0$, and the vectors $h'v - v$ span $\mathrm{span}(e_{5}, \dots, e_{8})$, so $M$ vanishes there. $M$ commutes with $j$, so on $\mathrm{span}(e_{3}, e_{4})$ it is $a + bj$. The reflection $\sigma$ reversing $e_{4}$, $e_{6}$, $e_{8}$ fixes $\mathrm{Re}\,\Omega$ and $e_{3}$ and sends $J$ to $-J$, so $U \mapsto \sigma U\sigma$ maps $SU(3)$ onto itself, preserves the Haar measure and $x$, and gives $M = \sigma M\sigma$, hence $b = 0$. So $M = c(\kappa)\Pi$ with $c = e_{3}\cdot M(e_{3}) = \int xe^{\kappa x}/\int e^{\kappa x}$, and the mean of $U^{-1}$ is $M^{\mathsf T} = c\Pi$. The check confirms $M = (\tfrac{\kappa}{6} - \tfrac{\kappa^{3}}{288})\Pi + O(\kappa^{5})$ entry by entry from the moments of Step 1: $c = \kappa m_{2} + \kappa^{3}(\tfrac16m_{4} - \tfrac12m_{2}^{2}) + O(\kappa^{5})$ with $m_{2} = \tfrac16$, $m_{4} = \tfrac1{16}$, and $\tfrac1{96} - \tfrac1{72} = -\tfrac1{288}$. Since $Ue_{3}$ is uniform on the sphere, $x$ and $-x$ have the same distribution, so the numerator of $c$ is $\int_{x > 0}x(e^{\kappa x} - e^{-\kappa x})\,dU > 0$ for $\kappa > 0$.

\emph{Step 4 (the frozen-triplet loop; proved).} For $C$ using $P$ distinct links once each, the links are independent after Step 2, and $\chi$ is linear, so $\langle\chi(U_{C})\rangle = \tfrac16\mathrm{tr}((c\Pi)^{P}) = \tfrac16c^{P}\mathrm{tr}\,\Pi = \tfrac13c^{P}$, using $\Pi^{P} = \Pi$ and $\mathrm{tr}\,\Pi = 2$ (checked for $P = 4, 6, 8$). At $\kappa = 0$, $c = 0$ and every such loop vanishes, as gauge-variant link averages do. For the weak doublet the stabiliser is trivial and the mean link is a scalar; for the triplet it is the projector onto one $J$-complex direction, the two others being acted on by the unbroken $SU(2)$. This is (ii).

\emph{Step 5 (Wilson quarks; proved).} For each flavour, the Grassmann integral of $\exp(-\bar\psi(1 - \kappa H)\psi)$ is $\det(1 - \kappa H)$ on spinor-colour space, with spinors the minimal left ideal $\Cl(4,0)f$, $f = \tfrac12(1 + d_{1}d_{2}d_{3}d_{4})$ (eight real dimensions), made complex by right multiplication by $d_{1}d_{2}$, which commutes with the left action and squares to $-1$, and colour the $J$-complex space $V$, the two complex units identified (the tensor product is over $\mathrm{span}\{1, j\}$). The real trace of left multiplication by $x$ on the ideal is $8\langle x\rangle_{0}$ (every basis blade other than $1$ anticommutes with some vector and is traceless), and for the products of hop factors used here the trace of left multiplication by $x$ composed with right multiplication by $d_{1}d_{2}$ vanishes (checked), so their complex spin trace is real and equals $4\langle x\rangle_{0}$. For small $\kappa$ on a finite lattice $\ln\det(1 - \kappa H) = -\sum_{n \ge 1}\tfrac{\kappa^{n}}{n}\mathrm{Tr}\,H^{n}$, and $\mathrm{Tr}\,H^{n}$ is the sum over closed walks of length $n$ of the spin trace times the colour trace of the products of hop factors; a walk that reverses a step gives $(1 - d_{\mu})(1 + d_{\mu}) = 1 - d_{\mu}^{2} = 0$. At $\beta = 0$, $\langle\chi(U_{C})\rangle = \int\chi(U_{C})\det(1 - \kappa H)^{n_{f}}\prod dU/\int\det(1 - \kappa H)^{n_{f}}\prod dU$. If a link of $C$ is traversed by none of the walks of a term, that term is linear in it and integrates to zero; so the first nonzero order is $\kappa^{P}$, from one closed walk of length $P$ covering the $P$ links of $C$, which is the boundary in one of its two orientations, counted from each of its $P$ starting points. Its contribution to $n_{f}\ln\det$ is $-n_{f}\kappa^{P}\bigl(s_{+}\mathrm{tr}_{J}(U_{C}) + s_{-}\mathrm{tr}_{J}(U_{C}^{-1})\bigr)$, with $s_{\pm}$ the spin traces of the two orientations and $\mathrm{tr}_{J}$ the $J$-complex colour trace, $\mathrm{tr}_{J}(U) + \mathrm{tr}_{J}(U^{-1}) = \mathrm{tr}\,U$. Along a side, $(1 \mp d_{\mu})^{2} = 2(1 \mp d_{\mu})$, so the product round an $R \times T$ rectangle is $2^{P - 4}(1 - d_{\mu})(1 - d_{\nu})(1 + d_{\mu})(1 + d_{\nu})$, whose scalar part is $2^{P - 4}(1 - 1 - 1 + (d_{\mu}d_{\nu})^{2}) = -2^{P - 3}$; so $s_{\pm} = -2^{P - 1}$ (the check multiplies the factors in $\Cl(4,0)$ for all six planes, $R, T \le 3$, both orientations). The order-$\kappa^{P}$ term of the weight is therefore $n_{f}2^{P - 1}\kappa^{P}\,\mathrm{tr}\,U_{C}$ (the check sums all closed walks of length $P$ on the $1 \times 1$ and $2 \times 1$ patches with links in $G$ and recovers $-2^{P - 1}\mathrm{tr}\,U_{C}$ for the walk sum), and the denominator is $1 + O(\kappa)$. With $U_{C}$ Haar distributed, $\int\chi(U)\,\mathrm{tr}\,U\,dU = \tfrac16\int(\mathrm{tr}\,U)^{2}dU = \tfrac26$ (degree two, the dimension of the invariants; also over $G$), so the leading term is $n_{f}2^{P - 1}\kappa^{P}\cdot\tfrac13 = \tfrac{n_{f}}{6}(2\kappa)^{P}$. This is (iii).

\emph{Step 6 (the comparison; argued).} At small $\beta > 0$ and small $\kappa > 0$ the leading terms of a large $R \times T$ loop are of order $(2\kappa)^{2(R + T)}$ (or $c^{2(R + T)}$) and $(\beta/18)^{RT}$, the leading area-law coefficient of the fundamental loop (from the gluing factor $\tfrac13$ and the weight $\beta/6$ per plaquette); the first dominates as the loop grows. This compares leading terms only and is not a bound on the full expectation. With triplets the fundamental Wilson loop is no order parameter of confinement, the colour analogue of \claimref{synthesis.non-abelian-gauge-fixing-has-gribov-copies-and-the-weak-rotor-lattice-with-doublets-has-a-perimeter-law}(iv); for scalars in the fundamental representation the absence of a phase boundary is Fradkin and Shenker's result (imported in the store's Fact on Singer's gauge slice, Neuberger's lattice BRST and Fradkin-Shenker complementarity, clause (iii), cited here in prose), and for quarks the same reading is argued, not proved. The series' colour sector has quark triplets, $n_{g}$ generations, an input of the stored one-loop running coefficients ($b_{3} = -11 + \tfrac43n_{g}$ there; $n_{g} = 3$ at the imported values); the Yang-Mills mass-gap problem is posed for the pure gauge field and is not addressed.

\emph{Status.} (i) computed exactly; (ii) and (iii) proved, with the expansions and walk sums checked exactly; the comparison argued. Standard lattice gauge theory recomputed in the colour block and in $\Cl(4,0)$; transport-absent.
